\documentclass[tikz,border=5mm]{standalone}
\usepackage{amsmath,amssymb}
\usepackage{fancyhdr}
\usetikzlibrary{calc,intersections}

\begin{document}
\begin{tikzpicture}[scale=1.2, font=\footnotesize]
    % Đường tròn ngoại tiếp (O)
    \coordinate (O) at (0,0);
    \def\R{3.4}
    \draw[thick, blue!80!black] (O) circle (\R);
    \filldraw[black] (O) circle (1.2pt) node[below right] {$O$};

    % Các đỉnh tam giác nhọn, không cân
    \coordinate (A) at (78:\R);
    \coordinate (B) at (205:\R);
    \coordinate (C) at (325:\R);

    % Chân các đường cao D, E, F
    \coordinate (D) at ($(B)!(A)!(C)$);
    \coordinate (E) at ($(A)!(B)!(C)$);
    \coordinate (F) at ($(A)!(C)!(B)$);

    % Trực tâm H
    \path[name path=altAD] (A) -- (D);
    \path[name path=altBE] (B) -- (E);
    \path[name intersections={of=altAD and altBE, by=H}];

    % Các trung điểm M, N, P
    \coordinate (M) at ($(A)!0.5!(B)$);
    \coordinate (N) at ($(A)!0.5!(C)$);
    \coordinate (P) at ($(B)!0.5!(C)$);

    % Điểm X = EF giao BC
    \path[name path=lineEF] ($(E)!-2.5!(F)$) -- ($(F)!3!(E)$);
    \path[name path=lineBC] ($(B)!-2.5!(C)$) -- ($(C)!1.5!(B)$);
    \path[name intersections={of=lineEF and lineBC, by=X}];

    % Điểm K = EF giao MN
    \path[name path=lineMN] ($(M)!-1.5!(N)$) -- ($(N)!3!(M)$);
    \path[name intersections={of=lineEF and lineMN, by=K}];

    % Tiếp tuyến tại A của (O) cắt MN tại Y
    % Vector tiếp tuyến vuông góc với OA
    \coordinate (tanA) at ($(A)!1!90:(O)$);
    \path[name path=lineTanA] ($(A)!-3!(tanA)$) -- ($(A)!3!(tanA)$);
    \path[name intersections={of=lineTanA and lineMN, by=Y}];

    % Đường thẳng Euler OH
    \draw[thick, red!80!black] ($(O)!-0.8!(H)$) -- ($(H)!2.2!(O)$);

    % Vẽ tam giác và các đoạn thẳng chính
    \draw[thick] (A) -- (B) -- (C) -- cycle;
    \draw[thick, blue] (X) -- (C);
    \draw[thick, blue] (X) -- (K);
    \draw[thick, purple] (Y) -- (N);
    \draw[thick, purple] (A) -- (Y);
    \draw[thick, orange!90!black] (A) -- (K);
    \draw[thick, orange!90!black] (X) -- (Y);

    % Đường trung tuyến AP và đường đối trung AL (cho câu b)
    \coordinate (G) at ($(A)!2/3!(P)$);
    % Dựng AL đối xứng AP qua phân giác góc A:
    % D là hình chiếu, L là trung điểm DP
    \coordinate (L) at ($(D)!0.5!(P)$);
    \draw[dashed, teal] (A) -- (P);
    \draw[dashed, olive] (A) -- (L);

    % Vẽ các đoạn đường cao
    \draw[gray, dashed] (A) -- (D);
    \draw[gray, dashed] (B) -- (E);
    \draw[gray, dashed] (C) -- (F);

    % Điểm
    \filldraw[black] (A) circle (1.5pt) node[above] {$A$};
    \filldraw[black] (B) circle (1.5pt) node[below left] {$B$};
    \filldraw[black] (C) circle (1.5pt) node[below right] {$C$};
    \filldraw[black] (D) circle (1.3pt) node[below] {$D$};
    \filldraw[black] (E) circle (1.3pt) node[above right] {$E$};
    \filldraw[black] (F) circle (1.3pt) node[above left] {$F$};
    \filldraw[black] (H) circle (1.3pt) node[below left] {$H$};
    \filldraw[black] (M) circle (1.3pt) node[left] {$M$};
    \filldraw[black] (N) circle (1.3pt) node[right] {$N$};
    \filldraw[black] (P) circle (1.3pt) node[below] {$P$};
    \filldraw[black] (X) circle (1.5pt) node[below left] {$X$};
    \filldraw[black] (K) circle (1.5pt) node[above] {$K$};
    \filldraw[black] (Y) circle (1.5pt) node[above left] {$Y$};
    \filldraw[black] (L) circle (1.3pt) node[below] {$L$};
    \filldraw[black] (G) circle (1.3pt) node[right] {$G$};

    % Đánh dấu góc vuông AK vuông góc OH
    \draw[red] ($(K)!0.15!(A)$) -- ++($(O)!0.15!(H)-(O)$);
\end{tikzpicture}
\end{document}
